%%% BEGIN fig:appP-dependency
% TEMPLATE of the generated App P graph (FIGURES_INDEX_v2 Sec. 4.9): node and edge data below are generated by
% test_ff/scripts/appP_template2.py from the \uses lists of FINAL_OUTLINE_v2 (a 19-statement slice); the integrator's
% generator replaces the slice by the selections of rule 1 and keeps the styles, the legend and the caption wording.
\begin{figure}[htbp]
\centering
\begin{tikzpicture}[fd,
  % node styles of FIGURES_INDEX_v2 Sec. 4.9 rule 3 (prefix ffP)
  ffP node/.style={draw,line width=0.45pt,rounded corners=1pt,font=\tiny\ttfamily,align=center,
                   inner xsep=1.5pt,inner ysep=1.1pt,minimum height=4mm},
  ffP R/.style={ffP node,fill=OliveGreen!15,draw=OliveGreen!60!black},
  ffP A/.style={ffP node,fill=RoyalBlue!15,draw=RoyalBlue!60!black},
  ffP P/.style={ffP node,fill=BurntOrange!15,draw=BurntOrange!70!black},
  ffP C/.style={ffP node,fill=Plum!15,draw=Plum!70!black},
  ffP V/.style={ffP node,fill=BrickRed!15,draw=BrickRed!70!black,double,double distance=0.6pt},
  ffP cond/.style={dash pattern=on 1.6pt off 1pt},
  ffP math/.style={ffP node,fill=white,draw=black!60},
  ffP def/.style={ffP node,fill=white,draw=black!45,line width=0.3pt,rounded corners=0pt},
  ffP conj/.style={ffP node,ellipse,inner xsep=0pt,inner ysep=0pt,fill=Plum!15,draw=Plum!70!black},
  ffP asm/.style={ffP node,fill=orange!8,draw=orange!65!black,dash pattern=on 2pt off 1.2pt},
  ffP unit/.style={font=\footnotesize,text=black!70},
  ffP sep/.style={draw=gray!40,line width=0.4pt,dash pattern=on 1pt off 1.5pt},
  ffP key/.style={font=\scriptsize,anchor=west,inner sep=0pt},
  % edge styles of rule 4
  ffP use/.style={draw=gray!60,line width=0.5pt,-{Stealth[length=2.6pt,width=2.2pt]}},
  ffP asmuse/.style={draw=orange!65!black,line width=0.5pt,dash pattern=on 2pt off 1.2pt,-{Stealth[length=2.6pt,width=2.2pt]}},
  ffP neck/.style={draw=red!60!black,line width=1.1pt,-{Stealth[length=3.6pt,width=3.2pt]}},
  ffP long/.style={opacity=0.5}]
  % ---- generated by test_ff/scripts/appP_template2.py (template of the App P generator) ----
  \node[ffP unit] at (0.000,2.847) {第~1~章};
  \node[ffP unit] at (2.501,2.847) {第~12~章};
  \node[ffP unit] at (5.747,2.847) {第~18~章};
  \node[ffP unit] at (8.583,2.847) {附录~N};
  \node[ffP unit] at (12.630,2.847) {附录~O};
  \draw[ffP sep] (0.517,2.677) -- (0.517,-2.476);
  \draw[ffP sep] (4.250,2.677) -- (4.250,-2.476);
  \draw[ffP sep] (7.152,2.677) -- (7.152,-2.476);
  \draw[ffP sep] (9.790,2.677) -- (9.790,-2.476);
  \node[ffP asm] (n00) at (0.000,0.375) {intro-\\bridge};
  \node[ffP A] (n01) at (1.593,0.425) {dy-factorization};
  \node[ffP conj] (n02) at (3.410,1.475) {dy-\\coloured-\\final};
  \node[ffP R,ffP cond] (n03) at (5.047,-1.675) {break-\\measurement};
  \node[ffP V] (n04) at (6.448,-1.775) {break-tmd};
  \node[ffP asm] (n05) at (7.949,-1.325) {appN-\\collins};
  \node[ffP asm] (n06) at (7.949,-0.325) {appN-\\cutsum};
  \node[ffP R,ffP cond] (n07) at (7.949,-2.125) {appN-\\subtraction};
  \node[ffP R,ffP cond] (n08) at (9.218,-1.125) {appN-dy};
  \node[ffP asm] (n09) at (10.531,-0.075) {appO-lcopt};
  \node[ffP def] (n10) at (10.531,1.375) {appO-\\family};
  \node[ffP asm] (n11) at (12.024,1.075) {appO-\\mixing};
  \node[ffP asm] (n12) at (12.024,0.175) {appO-\\deformation};
  \node[ffP asm] (n13) at (12.024,2.125) {appO-\\rewritten-cs};
  \node[ffP R] (n14) at (12.024,-0.675) {appO-\\target};
  \node[ffP R,ffP cond] (n15) at (13.461,0.075) {appO-main};
  \node[ffP R,ffP cond] (n16) at (14.730,-1.125) {appO-\\interface};
  \begin{scope}[on background layer]
  \draw[ffP use] (n00) -- (n01);
  \draw[ffP use] (n01) -- (n02);
  \draw[ffP use] (n03) -- (n04);
  \draw[ffP use] (n07) -- (n08);
  \draw[ffP asmuse] (n05) -- (n08);
  \draw[ffP asmuse] (n06) -- (n08);
  \draw[ffP use,ffP long] (n08) -- (n01);
  \draw[ffP asmuse,ffP long] (n09) -- (n03);
  \draw[ffP use] (n09) -- (n11);
  \draw[ffP use] (n09) -- (n14);
  \draw[ffP use] (n09) -- (n12);
  \draw[ffP use] (n09) -- (n13);
  \draw[ffP use] (n10) -- (n11);
  \draw[ffP use] (n10) -- (n12);
  \draw[ffP use] (n10) -- (n13);
  \draw[ffP asmuse] (n11) -- (n15);
  \draw[ffP use] (n14) -- (n15);
  \draw[ffP asmuse] (n12) -- (n15);
  \draw[ffP asmuse] (n13) -- (n15);
  \draw[ffP use] (n15) -- (n16);
  \draw[ffP use,ffP long] (n16) -- (n08);
  \draw[ffP neck] (n05) -- (n01);
  \draw[ffP neck] (n06) -- (n01);
  \draw[ffP neck] (n09) -- (n01);
  \draw[ffP neck] (n11) -- (n01);
  \draw[ffP neck] (n12) -- (n01);
  \draw[ffP neck] (n13) -- (n01);
  \end{scope}
  % ---------------------------------------------------------------- legend (rule 5)
  \begin{scope}[shift={(-0.397,-2.946)}]
    \node[ffP R,anchor=west,minimum width=4.5mm] (k1) at (0,0) {R};
    \node[ffP R,ffP cond,anchor=west,minimum width=4.5mm,right=0.8mm of k1] (k2) {R};
    \node[ffP A,anchor=west,minimum width=4.5mm,right=0.8mm of k2] (k3) {A};
    \node[ffP P,anchor=west,minimum width=4.5mm,right=0.8mm of k3] (k4) {P};
    \node[ffP C,anchor=west,minimum width=4.5mm,right=0.8mm of k4] (k5) {C};
    \node[ffP V,anchor=west,minimum width=4.5mm,right=0.8mm of k5] (k6) {V};
    \node[ffP key,right=1.2mm of k6] (t1) {严格性标签（虚线框为 $R_{\rm cond}$）};
    \node[ffP conj,anchor=west,minimum width=7mm,minimum height=3.4mm,right=3.5mm of t1] (k7) {};
    \node[ffP key,right=1mm of k7] (t2) {猜想};
    \node[ffP asm,anchor=west,minimum width=4.5mm,right=3.5mm of t2] (k8) {};
    \node[ffP key,right=1mm of k8] (t3) {编号假设};
    \node[ffP def,anchor=west,minimum width=4.5mm,right=3.5mm of t3] (k9) {};
    \node[ffP key,right=1mm of k9] (t4) {定义};
    \node[ffP math,anchor=west,minimum width=4.5mm,right=3.5mm of t4] (k10) {};
    \node[ffP key,right=1mm of k10] {纯数学};
    \draw[ffP use] (0,-0.5) -- ++(7mm,0) coordinate (e1);
    \node[ffP key] (t5) at ($(e1)+(1.2mm,0)$) {依赖（由前提指向陈述）};
    \draw[ffP asmuse] ($(t5.east)+(3.5mm,0)$) -- ++(7mm,0) coordinate (e2);
    \node[ffP key] (t6) at ($(e2)+(1.2mm,0)$) {编号假设 $\to R_{\rm cond}$ 陈述};
    \draw[ffP neck] ($(t6.east)+(3.5mm,0)$) -- ++(7mm,0) coordinate (e3);
    \node[ffP key] (t7) at ($(e3)+(1.2mm,0)$) {瓶颈依赖（主表“瓶颈”列）};
    \draw[ffP use,ffP long] ($(t7.east)+(3.5mm,0)$) -- ++(7mm,0) coordinate (e4);
    \node[ffP key] at ($(e4)+(1.2mm,0)$) {跨两列以上};
  \end{scope}
\end{tikzpicture}
\caption{本书规范陈述的依赖图（模板：Drell--Yan 与测量相容性部分的示意子图，17 个陈述；附录~P 的各图由各环境中的 \texttt{\string\uses} 数据按同一规则自动生成）。
各单元按书序自左向右分层，单元内按依赖深度分列；节点只印标签中类型前缀之后的部分（完整标签见附录~P 的状态表）；边由前提指向陈述，是可达关系在所画节点上的传递约化（经由未画出陈述的依赖也画成一条边）。
框的颜色为严格性标签，虚线框为 $R_{\rm cond}$ 条件定理，橙色虚线框为编号假设，椭圆为猜想，双线框为 [V]，方角细框为定义，白色圆角框为纯数学引理；
橙色虚线边为编号假设到 $R_{\rm cond}$ 陈述的依赖，粗红线为决定全阶论证地位的瓶颈依赖（取自主表的“瓶颈”列，即使可由其他边推出也照画），半透明线跨越两列以上。}
\label{fig:appP-dependency}
\end{figure}
%%% END fig:appP-dependency
